\section*{Bab \ref{ch:b1:cross-sections} --- Penampang geologi}

\begin{solution}{exo:b1:cross-sections:1}
Pada skala 1:25\,000, \qty{1}{cm} di peta sama dengan \qty{250}{m} di lapangan. Skala vertikalnya adalah \qty{1}{cm} untuk \qty{100}{m}, sehingga skala vertikal itu $250/100 = 2.5$ kali skala horizontal: $e = 2.5$.
\end{solution}

\begin{solution}{exo:b1:cross-sections:2}
$\gamma = \ang{90} - \ang{60} = \ang{30}$, dan $\tan\beta = \tan\ang{40}
\cos\ang{30} = 0.8391\times0.8660 = 0.7267$: $\beta = \ang{36.0}$, lapisannya turun ke arah ujung timur penampang. Pada penampang berazimut
\ang{000}, $\gamma = \ang{90}$ dan $\beta = 0$: penampang berjalan sepanjang \omterm{def:b1:geological-maps:attitude}{jurus} dan lapisan tampak horizontal.
\end{solution}

\begin{solution}{exo:b1:cross-sections:3}
$t = h\cos\delta = 120\times\cos\ang{12.16} = 120\times0.9776 =
\qty{117.3}{m}$.
\end{solution}

\begin{solution}{exo:b1:cross-sections:4}
$\tan\ang{12} = 0.2126$. Dengan $e = 2$: $\arctan 0.4251 = \ang{23.0}$; dengan
$e = 3$: $\arctan 0.6377 = \ang{32.5}$; dengan $e = 5$: $\arctan 1.0628 =
\ang{46.7}$. Sudut tumbuh jauh lebih lambat daripada $e$ itu sendiri, dan tidak pernah mencapai
\ang{90}.
\end{solution}

\begin{solution}{exo:b1:cross-sections:5}
(a) $t = 300\sin\ang{25} = \qty{126.8}{m}$. (b) $t = 300\sin\ang{35} =
\qty{172.1}{m}$. (c) $t = 300\sin\ang{15} = \qty{77.6}{m}$. Lebar \omterm{def:b1:geological-maps:geological-map}{singkapan} yang sama dapat berarti ketebalan yang berbeda lebih dari dua kali lipat: \omterm{def:b1:geological-maps:attitude}{kemiringan} tanah harus selalu diperhitungkan.
\end{solution}

\begin{solution}{exo:b1:cross-sections:6}
Dengan \omterm{def:b1:geological-maps:attitude}{arah kemiringan} $\alpha$ dan \omterm{def:b1:geological-maps:attitude}{kemiringan} $\delta$, teorema memberikan pada penampang utara--selatan ($\theta = 0$) $\tan\ang{20} = \tan\delta\cos\alpha$,
dan pada penampang timur--barat ($\theta = \ang{90}$) $\tan\ang{30} =
\tan\delta\cos(\alpha - \ang{90}) = \tan\delta\sin\alpha$. Dengan membagi,
$\tan\alpha = 0.5774/0.3640 = 1.586$, $\alpha = \ang{057.8}$; dengan menjumlahkan kuadratnya, $\tan\delta = \sqrt{0.3640^2 + 0.5774^2} = 0.6825$, $\delta =
\ang{34.3}$. Kedua \omterm{def:b1:cross-sections:apparent-dip}{kemiringan semu} lebih kecil daripada \omterm{def:b1:geological-maps:attitude}{kemiringan} sebenarnya, sebagaimana mestinya.
\end{solution}

\begin{solution}{exo:b1:cross-sections:7}
$100/\tan\ang{12.16} = 100/0.2155 = \qty{464.1}{m}$. Pada peta, \omterm{def:b1:cross-sections:structure-contour}{kontur struktur} sayap timur memang berjarak sedikit kurang dari setengah kilometer, berjarak sama dan sejajar: sayap itu adalah bidang.
\end{solution}

\begin{solution}{exo:b1:cross-sections:8}
B1: $524.02 - 488.98 = \qty{35.04}{m}$; B3: $529.98 - 494.93 =
\qty{35.05}{m}$. Di atas B2, dasar serpih adalah puncak batu gamping, \qty{180}{m} di atas dasarnya: $280.16 + 180 = \qty{460.16}{m}$,
lebih tinggi daripada tanah (\qty{380.64}{m}), itulah sebabnya ia tererosi di sana.
Antara B1 dan B2 \omterm{def:b1:geological-maps:geological-map}{kontak} itu naik sebesar $460.16 - 35.04 = \qty{425.12}{m}$
sepanjang \qty{2000}{m}: \omterm{def:b1:geological-maps:attitude}{kemiringan} $0.2126$, \omterm{def:b1:cross-sections:apparent-dip}{kemiringan semu} \ang{12.0}. B1
dan B3, pada jarak yang sama di kedua sisi B2, mendapati setiap \omterm{def:b1:geological-maps:geological-map}{kontak} pada elevasi yang sama dalam \qty{1}{cm}: \omterm{def:b1:geological-maps:fold}{lipatan} itu simetris terhadap bidang sumbu vertikal pada $x = \qty{3.5}{km}$, kedua sayapnya miring sama besar, dan tidak ada sesar yang memotong penampang di antara keduanya.
\end{solution}

\begin{solution}{exo:b1:cross-sections:9}
\omterm{def:b1:cross-sections:throw}{Lompatan horizontal} $= 150/\tan\ang{50} = 150/1.1918 = \qty{125.9}{m}$. Pada penampang berazimut \ang{150}, $\gamma = \ang{30}$ dan $\tan\beta = 1.1918\times0.8660 =
1.0321$: sesar tampak dengan \omterm{def:b1:cross-sections:apparent-dip}{kemiringan semu} \ang{45.9}. Pemotongan suatu lapisan horizontal berjarak \qty{150}{m} secara vertikal pada jejak itu, sehingga $150/1.0321 = \qty{145.3}{m}$ berjarak sepanjang penampang, lebih besar daripada \omterm{def:b1:cross-sections:throw}{lompatan horizontal}, karena penampang memotong sesar secara miring.
\end{solution}

\begin{solution}{exo:b1:cross-sections:10}
$\cos\gamma = \tan\ang{20}/\tan\ang{40} = 0.3640/0.8391 = 0.4338$, sehingga
$\gamma = \ang{64.3}$: penampang berazimut $\alpha \pm \ang{64.3}$ (dan arah yang berlawanan, tempat lapisan naik). \omterm{def:b1:cross-sections:apparent-dip}{Kemiringan semunya} adalah
\ang{40.0}, \ang{39.6}, \ang{38.3} dan \ang{36.0} untuk $\gamma = 0$, \ang{10},
\ang{20}, \ang{30}, tetapi \ang{22.8}, \ang{16.0}, \ang{8.3} dan $0$ untuk
$\gamma = \ang{60}$, \ang{70}, \ang{80}, \ang{90}. Memang, dengan menurunkan
$\tan\beta = \tan\delta\cos\gamma$,
$\dfrac{\mathrm d\beta}{\mathrm d\gamma} =
-\dfrac{\tan\delta\sin\gamma}{1 + \tan^2\delta\cos^2\gamma}$, yang nol
pada $\gamma = 0$ dan sama dengan $-\tan\delta$ pada $\gamma = \ang{90}$: penampang yang sedikit menyimpang dari \omterm{def:b1:geological-maps:attitude}{arah kemiringan} memperlihatkan hampir \omterm{def:b1:geological-maps:attitude}{kemiringan} sebenarnya, sementara dekat \omterm{def:b1:geological-maps:attitude}{jurus} setiap derajat \omterm{def:b1:spherical-astronomy:horizontal}{azimut} mengubah \omterm{def:b1:cross-sections:apparent-dip}{kemiringan semu} hampir satu derajat.
\end{solution}

\begin{solution}{exo:b1:cross-sections:11}
Lapisan miring ke timur, sehingga normal perlapisan, jika diikuti ke bawah, miring \ang{12.16} dari vertikal ke arah barat, ke arah \ang{260.7}.
(a) Lubang bor vertikal membentuk sudut $\psi = \ang{12.16}$ dengannya:
$176.0/\cos\ang{12.16} = \qty{180.0}{m}$, \omterm{def:b1:cross-sections:thickness}{ketebalan vertikal}. (b) Lubang bor yang menunjam \ang{60} ke arah \ang{260.7} miring \ang{30} dari vertikal pada sisi yang sama dengan normal: $\psi = \ang{30} - \ang{12.16} =
\ang{17.84}$, $\cos\psi = 0.9519$, dan inti sepanjang $176.0/0.9519 =
\qty{184.9}{m}$. (c) Ke arah \ang{080.7} ia miring pada sisi yang lain:
$\psi = \ang{30} + \ang{12.16} = \ang{42.16}$, $\cos\psi = 0.7413$, dan
inti sepanjang $176.0/0.7413 = \qty{237.4}{m}$. Inti paling pendek, sama dengan
$t$, bila lubang bor berjalan sepanjang normal, menunjam \ang{77.84} ke arah
\ang{260.7}.
\end{solution}

\begin{solution}{exo:b1:cross-sections:12}
Pada bidang vertikal timur--barat, dengan sudut diukur dari timur, lapisan atas naik ke arah timur sebesar \ang{20} dan lapisan bawah turun ke arah timur sebesar \ang{70}. Memutar \omterm{def:b1:geological-maps:geological-map}{singkapan} terhadap \omterm{def:b1:geological-maps:attitude}{jurusnya} sebesar \ang{20}, dengan sisi timur ke bawah, membuat lapisan atas mendatar dan membawa lapisan bawah ke $\ang{70} + \ang{20} = \ang{90}$: lapisan itu tegak ketika lapisan atas diendapkan. Sejarahnya, dari peristiwa tertua: lapisan bawah diendapkan datar; ia dimiringkan hingga tegak, lalu diangkat; tepi-tepinya tererosi menjadi suatu permukaan; lapisan atas diendapkan melintasi permukaan itu; seluruhnya dimiringkan \ang{20} turun ke barat; erosi memotong permukaan tanah hari ini. Itulah sejarah \omterm{def:b1:geological-maps:geological-map}{singkapan} Hutton di Siccar Point.
\end{solution}

\begin{solution}{pb:b1:cross-sections:1}
\textbf{1.} Ketebalan batuan di atas terowongan adalah $z - 200$: terbesar pada
$x = \qty{1}{km}$, $556 - 200 = \qty{356}{m}$; terkecil pada $x =
\qty{9}{km}$, $339 - 200 = \qty{139}{m}$.

\textbf{2.} \qty{1}{km} elevasi berukuran \qty{4.2}{cm} dan \qty{1}{km} jarak berukuran \qty{1.4}{cm}: $e = 4.2/1.4 = 3$.

\textbf{3.} $380.64 - 220.48 = \qty{160.16}{m}$.

\textbf{4.} Di B2 dasar batu gamping berada pada $380.64 - 100.48 =
\qty{280.16}{m}$ dan dasar batu pasir pada \qty{160.16}{m}: terowongan, pada \qty{200}{m}, berada di batu pasir.

\textbf{5.} B1: $524.02 - 668.98 = \qty{-144.96}{m}$; B3: $529.98 - 674.93 =
\qty{-144.95}{m}$.

\textbf{6.} Dari barat ke timur: kapur, serpih, batu gamping, batu pasir (di inti \omterm{def:b1:geological-maps:fold}{antiklin}), batu gamping, serpih, dan kapur lagi sampai $x = \qty{9}{km}$, melewati \omterm{def:b1:geological-maps:fold}{sinklin} dan memotong sesar.

\textbf{7.} $\gamma = \ang{80.7} - \ang{90} = \ang{-9.3}$; $\tan\beta =
\tan\ang{12.16}\cos\ang{9.3} = 0.2155\times0.9869 = 0.2126$: $\beta =
\ang{12.0}$, turun ke arah timur.

\textbf{8.} $(280.16 + 144.95)/2000 = 425.11/2000 = 0.2126$: \ang{12.0}.

\textbf{9.} $\arctan(3\times0.2126) = \arctan 0.6377 = \ang{32.5}$.

\textbf{10.} $\gamma = \ang{120} - \ang{90} = \ang{30}$; $\tan\beta =
\tan\ang{60}\cos\ang{30} = 1.732\times0.866 = 1.500$: $\beta = \ang{56.3}$.

\textbf{11.} Sayap-sayapnya miring ke arah \ang{080.7} atau \ang{279.3}, hanya \ang{9.3} dari garis penampang, dan $\cos\ang{9.3} = 0.987$; sesar miring ke arah \ang{120}, \ang{30} darinya, dan $\cos\ang{30} = 0.866$.

\textbf{12.} Bidang sesar adalah $z = 250 - \tan\ang{60}\,[(x - 4300)
\sin\ang{120} + (y - 1000)\cos\ang{120}]$. Dengan $y = 6000$, $(y - 1000)
\cos\ang{120} = -2500$, sehingga $(x - 4300)\sin\ang{120} = 2500 - (z - 250)/
\tan\ang{60}$, dan $1/\tan\ang{60} = \tan\ang{30}$.

\textbf{13.} Pada $z = 200$: $x = 4300 + (2500 + 50\times0.5774)/0.8660 =
4300 + 2528.9/0.8660 = \qty{7220}{m}$, yaitu \qty{7.22}{km} dari C.

\textbf{14.} \omterm{def:b1:cross-sections:throw}{Lompatan horizontal} $= 200/\tan\ang{60} = \qty{115.5}{m}$; pemisahan sepanjang C--D $= 200/1.500 = \qty{133.3}{m}$.

\textbf{15.} Sesar miring ke timur-tenggara dan merupakan sesar normal: dinding gantung, di timur, turun. Pada tingkat terowongan, kapur berada di kedua sisi sesar: di barat sesar, kapur \omterm{def:b1:geological-maps:fold}{sinklin}; di timur sesar, kapur yang diturunkan \qty{200}{m}.

\textbf{16.} Terowongan tetap berada di kapur melintasi sesar, sehingga tidak ada perubahan batuan di muka kerja yang memberi tanda; hanya zona batuan yang remuk, dan mungkin air, yang melakukannya. Para insinyur memperkirakan posisinya dari penampang (di sini \qty{7.22}{km}), mengebor lubang uji di depan mesin ketika ia mendekat, dan menyiapkan penyangga untuk zona yang remuk.

\textbf{17.} $668.98 - 488.98 = \qty{180.00}{m}$.

\textbf{18.} $t = 180\cos\ang{12.16} = \qty{176.0}{m}$.

\textbf{19.} Di puncak ($x = 3500$) dasar batu pasir berada pada \qty{160.16}{m} (soal 3); ia turun di kedua sayap dengan \omterm{def:b1:cross-sections:apparent-dip}{kemiringan semu} \ang{12.0} dari soal 7, yaitu sebesar $\tan\ang{12.0}$ per meter jarak horizontal dari sumbu.

\textbf{20.} Dasar batu gamping berada \qty{120}{m} di atas dasar batu pasir, pada $280.16 - u\tan\ang{12.0}$ dengan $u = |x - 3500|$, dan puncaknya \qty{180}{m} lebih tinggi, pada $460.16 - u\tan\ang{12.0}$. Dengan $\tan\ang{12.0} =
0.21256$: dasarnya berada di bawah \qty{200}{m} bila $u > 80.16/0.21256 = 377.1$, dan puncaknya di atas bila $u < 260.16/0.21256 = 1224.0$.

\textbf{21.} $x = 3500 \pm 377.1$ dan $3500 \pm 1224.0$: terowongan berada di batu gamping dari \qty{2276}{m} sampai \qty{3123}{m} dan dari \qty{3877}{m} sampai \qty{4724}{m}.

\textbf{22.} Sepanjang garis horizontal, dasar dan puncak batu gamping, yang berjarak \qty{180}{m} secara vertikal dan keduanya miring dengan $\tan\beta$, dilintasi pada jarak horizontal $h/\tan\beta$ satu sama lain: $180/0.21256 = \qty{846.8}{m}$ pada setiap sayap, sama dengan $1224.0 - 377.1 = 846.9$ dalam pembulatan.

\textbf{23.} $2\times846.8 = \qty{1693.7}{m}$: terowongan menembus
\textbf{sekitar \qty[text-series-to-math=true]{1694}{m} batu gamping}, lebih dari seperlima panjang \qty{8}{km}-nya, dalam dua bentangan sepanjang \qty{847}{m} di kedua sisi inti batu pasir \omterm{def:b1:geological-maps:fold}{antiklin}.
\end{solution}
